\chapter{Preliminary Analysis}
\label{chapter-PreliminaryAnalysis}



\section{Preliminary Design Step I: Design inputs from mission and failure analysis}
* Mi servono degli input: prendo le informazioni raccolte nella fase preliminare di analisi missione (ref analisi missione)per stabilire probabilità di impatto and what not. 

* traiettoria
* punto di partenza, punto di arrivo
* drone scelto
* payload scelto
* medicinali scelti per simulazione

* serve energia di impatto?

---> risultati del WP2


\section{Preliminary Design Step II: Payload container allowables}

* impact of container in LS-Dyna and what not. 
* diversi angoli di assetto per vedere il caso peggiore di accelerazioni locali
* cosa  possono sopportare le fialette e il contenitore????
* tipi di impatto, contenitore tipico che le contiene (menzionare la tesi del tipello) e che per il dimensionamento non viene considerato, approssimato a infinitamente rigido e che trasmette le accelerazioni di impatto così come le riceve ----> airbag deve essere dimensionato così
--------> tutte le analisi statistiche di prima, che hanno tenuto in considerazione anche il fatto che prima o poi si sorvolerà la città una volta risolti i cavilli burocratici che impongono una certa mitigazione del rischio ---> airbag. 



----> risultati del WP3



----> CONCLUSIONE: I have what I need to design the damn airbag. 

